% Part: incompleteness
% Chapter: incompleteness-provability
% Section: lob-thm

\documentclass[../../../include/open-logic-section]{subfiles}

\begin{document}

\olfileid{inc}{inp}{tar}

\olsection{సత్యపు నిర్వచనాతీతత}

\emph{నిర్వచనీయత} అనే భావన అంకగణిత భాషకు అధికారిక అర్థవిజ్ఞానం ఉండడంపై
ఆధారపడుతుంది. అంకగణిత భాషలోని సూత్రాలు, వాక్యాల సమితిని మనం వివరించాం.
అలాంటి వాక్యాలు సహజ సంఖ్యల గురించి ప్రకటనలు చేస్తాయని చదవడమే ``ఉద్దేశిత
అర్థనిర్దేశం''; అలాంటి ప్రకటన సత్యమో అసత్యమో కావచ్చు. అంకగణిత భాషలోని
సంకేతాలకు ప్రమాణ అర్థనిర్దేశమూ, వ్యక్తి క్షేత్రంగా $\Nat$యూ గల
\tetoken{నిర్మాణం}{!!{structure}} $\Struct{N}$గా ఉండనివ్వండి. అప్పుడు $\Sat{N}{!A}$ అంటే
``ప్రమాణ అర్థనిర్దేశంలో $!A$ సత్యం'' అని అర్థం.

\begin{defn}
ప్రతి $n_1,\dots,n_k$కు, $R(n_1,\dots,n_k)$ అవడం, $\Sat{N}{!A(\num
n_1,\dots,\num n_k)}$ అవడం సమానార్థకమయ్యే అంకగణిత భాషలోని సూత్రం
$!A(x_1,\dots,x_k)$ ఉంటే, అప్పుడు మాత్రమే సహజ సంఖ్యల సంబంధం
$R(x_1,\dots,x_k)$ను $\Struct{N}$లో \emph{నిర్వచించదగినది} అంటాం.
\end{defn}

మరో విధంగా చెప్పాలంటే, $\Th{TA} = \Setabs{!A}{\Sat{N}{!A}}$ అంకగణితపు
సత్య వాక్యాల సమితి అయినప్పుడు, సంబంధం $\Th{TA}$ సిద్ధాంతంలో ప్రాతినిధ్యం
చేయదగినదైతే, అప్పుడు మాత్రమే అది $\Struct{N}$లో నిర్వచించదగినది. (ఇది
వెంటనే స్పష్టంగా లేకపోతే, నిర్వచనాలను తిరిగి చూసి ఇది నిజమేనని మిమ్మల్ని
మీరు ఒప్పించుకోవాలి.)

\begin{lem}
ప్రతి గణనీయ సంబంధమూ~$\Struct{N}$లో నిర్వచించదగినది.
\end{lem}

\begin{proof}
$\Th{Q}$లో ఒక సంబంధానికి ప్రాతినిధ్యం వహించే సూత్రమే $\Struct{N}$లో కూడా
అదే సంబంధాన్ని నిర్వచిస్తుందని సులభంగా తనిఖీ చేయవచ్చు.
\end{proof}

ఇప్పుడు దీని విపర్యయం కూడా నిజమేనా అని అడగవచ్చు. అంటే $\Struct{N}$లో
నిర్వచించదగిన ప్రతి సంబంధమూ గణనీయమైనదేనా? సమాధానం కాదు. ఉదాహరణకు:

\begin{lem}
నిలుపుదల సంబంధం $\Struct{N}$లో నిర్వచించదగినది.
\end{lem}

\begin{proof}
ప్రతి పాక్షిక గణనీయ ప్రమేయం~$f$కు, అన్ని $x \in \Nat$లకు $f(x) =
U(\umin{s}{T(e,x,s)})$ అయ్యే సూచిక~$e$ ఉంటుందని క్లీనీ సాధారణ రూప
సిద్ధాంతం చెబుతుందని గుర్తుచేసుకోండి; ఇక్కడ $U$, $T$ ఆదిమ పునరావర్తన
ప్రమేయాలు, అందువల్ల సర్వత్ర నిర్వచితమైనవి. కాబట్టి $T(e,x,s)$ అయ్యే ఒక~$s$
ఉంటే, అప్పుడు మాత్రమే $f(x)$ నిర్వచితమవుతుంది (అంటే గణన ఆగుతుంది).

ఇప్పుడు $H$ నిలుపుదల సంబంధం అనుకుందాం, అంటే
\[
H = \Setabs{\tuple{e,x}}{\lexists[s][T(e, x, s)]}.
\]
$!D_T$ అనేది $\Struct{N}$లో $T$ను నిర్వచించనివ్వండి. అప్పుడు
\[
H = \Setabs{\tuple{e,x}}{\Sat{N}{\lexists[s][!D_T(\num e, \num x, s)]}},
\]
కాబట్టి $\lexists[s][!D_T(z, x, s)]$ అనేది $\Struct{N}$లో~$H$ను
నిర్వచిస్తుంది.
\end{proof}

\begin{prob}
$Q(n) \defiff n \in \Setabs{\Gn{!A}}{\Th{Q} \Proves !A}$ అంకగణితంలో
నిర్వచించదగినదని చూపండి.
\end{prob}

$\Th{TA}$ సంగతేమిటి? అది అంకగణితంలో నిర్వచించదగినదా? అంటే,
$\Setabs{\Gn{!A}}{\Sat{N}{!A}}$ సమితి అంకగణితంలో నిర్వచించదగినదా?
టార్స్కీ సిద్ధాంతం దీనికి ప్రతికూల సమాధానం ఇస్తుంది.

\begin{thm}
\ollabel{thm:tarski}
అంకగణితపు సత్య \tetoken{వాక్యాల}{!!{sentence}s} సమితి అంకగణితంలో నిర్వచించదగినది కాదు.
\end{thm}

\begin{proof}
$!D(x)$ దానిని నిర్వచిస్తుందని అనుకుందాం, అంటే $\Sat{N}{!A}$ అవడం,
$\Sat{N}{!D(\gn{!A})}$ అవడం సమానార్థకం. స్థిరబిందు ఉపసిద్ధాంతం ప్రకారం,
$\Th{Q} \Proves !A \liff \lnot !D(\gn{!A})$ అయ్యే వాక్యం $!A$ ఉంటుంది;
అందువల్ల $\Sat{N}{!A \liff \lnot !D(\gn{!A})}$. కానీ అప్పుడు
$\Sat{N}{!A}$ అవడం, $\Sat{N}{\lnot !D(\gn{!A})}$ అవడం సమానార్థకం;
ఇది $!D(y)$ అంకగణితపు సత్య ప్రకటనల సమితిని నిర్వచిస్తుందనే ఊహకు
విరుద్ధం.
\end{proof}

టార్స్కీ ఈ విశ్లేషణను సత్యమనే మరింత సాధారణ తాత్త్విక భావనకు అన్వయించాడు.
ఏ భాష $L$ ఇచ్చినా, $L$కు సరిపడే సత్యభావన ప్రతి వాక్యం $X$కు కింది నియమాన్ని
తృప్తిపరచాలని టార్స్కీ వాదించాడు:
\begin{quote}
`$X$' సత్యమవడం, $X$ అవడం సమానార్థకం.
\end{quote}
ఆంగ్లానికి టార్స్కీ తరచూ ఉటంకించబడే ఉదాహరణ ఈ వాక్యం:
\begin{quote}
`మంచు తెల్లగా ఉంటుంది' అనేది సత్యమవడం, మంచు తెల్లగా ఉండడం సమానార్థకం.
\end{quote}
అయితే వికర్ణ ప్రమేయానికి ప్రాతినిధ్యం వహించగలిగేంత బలమైన ఏ భాషలోనైనా,
ఏ భాషా విధేయం $T(x)$కైనా, ``$X$ అవడం, $T(\text{`$X$'})$ కాకపోవడం
సమానార్థకం'' అనే షరతును తృప్తిపరచే వాక్యం $X$ను మనం నిర్మించవచ్చు. ఒక సత్య
విధేయం కొన్ని వాక్యాలను సత్యమనీ అసత్యమనీ రెండూ ప్రకటించకూడదని మనం కోరుతాం
కాబట్టి, భాష హద్దుల వెలుపలికి ఏదో విధంగా వెళ్లకుండా ఆ భాషలోని అన్ని
వాక్యాలకు సత్య విధేయాన్ని నిర్దేశించలేమని టార్స్కీ నిర్ధారించాడు. మరో విధంగా,
ఒక భాషకు చెందిన సత్య విధేయాన్ని అదే భాషలో నిర్వచించలేం.

\end{document}
